% GENERATED FILE. Edit canonical lessons, then run npm run generate:books.
% canonical-source-hash: fnv1a64:a8a02c4ae3684a9a
% canonical-lessons: FR-C106-velo, FR-C106-metro, FR-C106-taxi, FR-C106-parking, FR-C106-ascenseur

\chapter{Bicycle, Underground, Taxi, Car park, Lift}
\label{ch:fr-bicycle-underground-taxi-car-park-lift}
\hlchaptermodality{106}

\begin{chapteropening}
\textbf{By the end of this chapter:} \emph{I can name a bicycle, the underground, a taxi, a car park, a lift.}

Complete the last lesson of chapter 106: I can name a bicycle, the underground, a taxi, a car park, a lift.
\end{chapteropening}

\section[le vélo]{le vélo — a bicycle}

\label{lesson:FR-C106-velo}



\begin{quote}
Before the new one: say the French for fifty, then the French for a hundred.
\end{quote}

\subsection*{You'll want to know: le vélo}
\textbf{le vélo} — \textquotedblleft{}a bicycle\textquotedblright{}.

\begin{cousinweb}
\textbf{vélo} — short for \emph{vélocipède}, a 19th-century coinage from Latin \emph{velox}, \textquotedblleft{}swift\textquotedblright{}, and \emph{pes}, \emph{pedis}, \textquotedblleft{}foot\textquotedblright{} (EtymDB 2.1 gives the two Latin parts).
\end{cousinweb}

\begin{grammarlens}[title={What you've built}]
One more word for this chapter.
\end{grammarlens}

\subsection*{Your turn}
\begin{itemize}
\raggedright
  \item \emph{Say it:} \emph{le vélo}
  \item \emph{Say it:} \emph{le vélo}, once more
  \item \emph{Say it:} say \emph{le vélo}
  \item \emph{Recall:} say the French for fifty, then the French for a hundred, then say \emph{le vélo} again
\end{itemize}

\begin{tcolorbox}[breakable,colback=teal!4,colframe=teal!35!black,title={Before you move on}]
What does le vélo mean? (\textquotedblleft{}A bicycle\textquotedblright{}.) Say it once more.
\end{tcolorbox}
\section[le métro]{le métro — the underground, the metro}

\label{lesson:FR-C106-metro}



\begin{quote}
Before the new one: say the French for a hundred, then the French for a bicycle.
\end{quote}

\subsection*{You'll want to know: le métro}
\textbf{le métro} — \textquotedblleft{}the underground, the metro\textquotedblright{}.

\begin{cousinweb}
\textbf{métro} — short for \emph{métropolitain}, the city railway, from Greek \emph{metropolis}, \textquotedblleft{}a mother city\textquotedblright{}, built from \emph{meter}, \textquotedblleft{}mother\textquotedblright{}, and \emph{polis}, \textquotedblleft{}city\textquotedblright{}. (EtymDB 2.1)
\end{cousinweb}

\begin{grammarlens}[title={What you've built}]
One more word for this chapter.
\end{grammarlens}

\subsection*{Your turn}
\begin{itemize}
\raggedright
  \item \emph{Say it:} \emph{le métro}
  \item \emph{Say it:} \emph{le métro}, once more
  \item \emph{Say it:} say \emph{le métro}
  \item \emph{Recall:} say the French for a hundred, then the French for a bicycle, then say \emph{le métro} again
\end{itemize}

\begin{tcolorbox}[breakable,colback=teal!4,colframe=teal!35!black,title={Before you move on}]
What does le métro mean? (\textquotedblleft{}The underground, the metro\textquotedblright{}.) Say it once more.
\end{tcolorbox}
\section[le taxi]{le taxi — a taxi}

\label{lesson:FR-C106-taxi}



\begin{quote}
Before the new one: say the French for a bicycle, then the French for the underground.

Ask for a taxi: \emph{je demande un taxi}. \textbf{Demander} is to ask, never to demand.
\end{quote}

\subsection*{You'll want to know: le taxi}
\textbf{le taxi} — \textquotedblleft{}a taxi\textquotedblright{}.

\begin{cousinweb}
\textbf{taxi} — short for \emph{taximètre}, the fare meter, built from Medieval Latin \emph{taxa}, \textquotedblleft{}a charge\textquotedblright{}, and Greek \emph{metron}, \textquotedblleft{}measure\textquotedblright{}.
\end{cousinweb}

\begin{grammarlens}[title={What you've built}]
One more word for this chapter.
\end{grammarlens}

\subsection*{Your turn}
\begin{itemize}
\raggedright
  \item \emph{Say it:} \emph{le taxi}
  \item \emph{Say it:} \emph{le taxi}, once more
  \item \emph{Say it:} say \emph{le taxi}
  \item \emph{Recall:} say the French for a bicycle, then the French for the underground, then say \emph{le taxi} again
\end{itemize}

\begin{tcolorbox}[breakable,colback=teal!4,colframe=teal!35!black,title={Before you move on}]
What does le taxi mean? (\textquotedblleft{}A taxi\textquotedblright{}.) Say it once more.
\end{tcolorbox}
\section[le parking]{le parking — a car park}

\label{lesson:FR-C106-parking}



\begin{quote}
Before the new one: say the French for the underground, then the French for a taxi.
\end{quote}

\subsection*{You'll want to know: le parking}
\textbf{le parking} — \textquotedblleft{}a car park\textquotedblright{}.

\begin{cousinweb}
\textbf{parking} $\leftarrow$ English \emph{parking} — borrowed in the 20th century; English itself says \emph{car park}.
\end{cousinweb}

\begin{grammarlens}[title={What you've built}]
One more word for this chapter.
\end{grammarlens}

\subsection*{Your turn}
\begin{itemize}
\raggedright
  \item \emph{Say it:} \emph{le parking}
  \item \emph{Say it:} \emph{le parking}, once more
  \item \emph{Say it:} say \emph{le parking}
  \item \emph{Recall:} say the French for the underground, then the French for a taxi, then say \emph{le parking} again
\end{itemize}

\begin{tcolorbox}[breakable,colback=teal!4,colframe=teal!35!black,title={Before you move on}]
What does le parking mean? (\textquotedblleft{}A car park\textquotedblright{}.) Say it once more.
\end{tcolorbox}
\section[l'ascenseur]{l'ascenseur — a lift, an elevator}

\label{lesson:FR-C106-ascenseur}



\begin{quote}
Before the new one: say the French for a taxi, then the French for a car park.
\end{quote}

\subsection*{You'll want to know: l'ascenseur}
\textbf{l'ascenseur} — \textquotedblleft{}a lift, an elevator\textquotedblright{}.

\begin{cousinweb}
\textbf{ascenseur} $\leftarrow$ Latin \emph{ascensor}, \textquotedblleft{}one who climbs\textquotedblright{}, from \emph{ascendere}, \textquotedblleft{}to climb up\textquotedblright{} — a learned borrowing. (EtymDB 2.1)
\end{cousinweb}

\begin{grammarlens}[title={What you've built}]
One more word for this chapter.
\end{grammarlens}

\subsection*{Your turn}
\begin{itemize}
\raggedright
  \item \emph{Say it:} \emph{l'ascenseur}
  \item \emph{Say it:} \emph{l'ascenseur}, once more
  \item \emph{Say it:} say \emph{l'ascenseur}
  \item \emph{Recall:} say the French for a taxi, then the French for a car park, then say \emph{l'ascenseur} again
\end{itemize}

\begin{tcolorbox}[breakable,colback=teal!4,colframe=teal!35!black,title={Before you move on}]
What does l'ascenseur mean? (\textquotedblleft{}A lift, an elevator\textquotedblright{}.) Say it once more.
\end{tcolorbox}
